\chapter{1977年安徽卷（理）}

\section{解答题}

\begin{problem}
\begin{enumerate}
    \item 计算：\(0.25\times\left( -2 \right)^2-4\div\left( \sqrt{5}-1 \right)^0-\left( \frac{1}{6} \right)^{-\frac{1}{2}}+\frac{\sqrt{3}}{\sqrt{3}-\sqrt{2}}\).
    \item 求函数 \(y=\log_x\left( 4^x-16 \right)\) 的定义域.
    \item 求 \(\sin\left( -915^\circ \right)\) 的值.
    \item 求函数 \(y=\arcsin\left( 1-2x \right)\) 的定义域.
    \item 化复数 \(-1+\i\) 为三角函数式.
    \item 化简：\(\sec A\cdot\sqrt{1-\sin^2 A}\)（\(90^\circ<A\leqslant180^\circ\)）.
    \item 已知圆锥的母线长为 \(3 \mathrm{~cm}\)，底面直径 \(2 \mathrm{~cm}\)，求它的侧面积（\(\pi\) 取 \(3.14\)）.
\end{enumerate}
\end{problem}

\begin{solution}
\begin{enumerate}
\item 由于 \(\sqrt{5}-1\neq0\)，有 \(\left(\sqrt{5}-1\right)^0=1\)，且 \(\left(\frac{1}{6}\right)^{-\frac{1}{2}}=\sqrt{6}\). 对最后一项分母有理化，得
\[
\begin{gathered}
0.25\times\left(-2\right)^2-4\div\left(\sqrt{5}-1\right)^0-\left(\frac{1}{6}\right)^{-\frac{1}{2}}+\frac{\sqrt{3}}{\sqrt{3}-\sqrt{2}} \\
=\frac{1}{4}\times4-4-\sqrt{6}+\sqrt{3}\left(\sqrt{3}+\sqrt{2}\right)
=1-4-\sqrt{6}+3+\sqrt{6}=0
\end{gathered}
\]

\item 对数的底数须满足 \(x>0\) 且 \(x\neq1\)，真数须满足 \(4^x-16>0\). 因为 \(4^x\) 是增函数，\(4^x>16=4^2\) 等价于 \(x>2\). 因而不等式组
\[
\begin{cases}
x>0,\\
x\neq1,\\
4^x-16>0
\end{cases}
\]
的解为 \(x>2\). 所以函数 \(y=\log_x\left(4^x-16\right)\) 的定义域为 \(\left(2,+\infty\right)\).

\item 利用正弦函数的奇偶性和周期性，
\[
\begin{gathered}
\sin\left(-915^\circ\right)=-\sin915^\circ=-\sin195^\circ=\sin15^\circ \\
=\sin\left(45^\circ-30^\circ\right)
=\sin45^\circ\cos30^\circ-\cos45^\circ\sin30^\circ \\
=\frac{\sqrt{2}}{2}\cdot\frac{\sqrt{3}}{2}-\frac{\sqrt{2}}{2}\cdot\frac{1}{2}
=\frac{\sqrt{6}-\sqrt{2}}{4}.
\end{gathered}
\]

\item 反正弦函数的自变量必须属于 \(\left[-1,1\right]\)，所以应有
\[
-1\leqslant1-2x\leqslant1
\]
由 \(-1\leqslant1-2x\) 得 \(-2\leqslant-2x\)，即 \(x\leqslant1\)；由 \(1-2x\leqslant1\) 得 \(-2x\leqslant0\)，即 \(x\geqslant0\). 两者取交集，故定义域为 \(\left[0,1\right]\).

\item 复数 \(-1+\i\) 的模为 \(\sqrt{\left(-1\right)^2+1^2}=\sqrt{2}\). 它的实部为负、虚部为正，辐角在第二象限，且 \(\cos\theta=-\frac{\sqrt{2}}{2}\)，\(\sin\theta=\frac{\sqrt{2}}{2}\)，因此可取 \(\theta=\frac{3\pi}{4}\). 所以
\[
-1+\i=\sqrt{2}\left(\cos\frac{3\pi}{4}+\i\sin\frac{3\pi}{4}\right)
\]

\item 因为 \(90^\circ<A\leqslant180^\circ\)，所以 \(\cos A<0\). 由平方关系和平方根的非负性，
\[
\sec A\cdot\sqrt{1-\sin^2 A}
=\sec A\cdot\sqrt{\cos^2 A}
=\sec A\cdot\left|\cos A\right|
=\sec A\cdot\left(-\cos A\right)=-1
\]

\item 底面半径为 \(R=\frac{2}{2}\mathrm{~cm}=1\mathrm{~cm}\)，母线长为 \(L=3\mathrm{~cm}\). 圆锥侧面展开后是半径为 \(L\) 的扇形，其弧长等于底面周长 \(C=2\pi R\)，所以侧面积为
\[
S_{\text{侧}}=\frac{1}{2}CL=\frac{1}{2}\times2\pi R\times L=\pi RL=3.14\times1\times3=9.42\mathrm{~cm}^2
\]
\end{enumerate}
\end{solution}

\begin{problem}
用宽为 \(80 \mathrm{~cm}\) 的铁皮制作一个槽，槽的截面 \(ABCD\) 是一个矩形：

\begin{enumerate}
    \item 求槽的截面积 \(y\) 关于槽深 \(x\) 的函数关系式.
    \item 求出这个函数图象的顶点坐标，对称轴方程.
    \item 当槽深为多少厘米时，槽的截面积最大？这个最大面积是多少平方厘米？
\end{enumerate}

\bitmapfigure[width=5.6cm]{img/1977/anhui_science/q2_fig1.png}
\end{problem}

\begin{solution}
\begin{enumerate}
\item 铁皮总宽为 \(80\mathrm{~cm}\)，两侧各折出宽为 \(x\) 的边，因此槽底宽为 \(80-2x\). 槽深和槽底宽都应为正数，故 \(0<x<40\). 截面积为
\(y=x\left(80-2x\right)=-2x^2+80x\)，\(0<x<40\)

\item 配方得
\[
y=-2x^2+80x=-2\left(x-20\right)^2+800
\]
所以函数图象的顶点坐标为 \(\left(20,800\right)\)，对称轴方程为 \(x=20\). 因为二次项系数为负，图象开口向下.

\item 在定义域 \(0<x<40\) 内，\(\left(x-20\right)^2\geqslant0\)，所以 \(y=-2\left(x-20\right)^2+800\leqslant800\)，等号当且仅当 \(x=20\). 因此槽深为 \(20\mathrm{~cm}\) 时截面积最大，最大截面积为 \(800\mathrm{~cm}^2\).
\end{enumerate}
\end{solution}

\begin{problem}
\begin{enumerate}
    \item 写出 \(\left( \sqrt[3]{x}-\frac{2}{3\sqrt[5]{x}} \right)^n\) 的展开式的第 \(\left( r+1 \right)\) 项.
    \item 设 \(AM\) 是直角三角形斜边 \(BC\) 上的中线，求证：\(BC^2+AC^2+AB^2=8AM^2\).
\end{enumerate}

\bitmapfigure[width=3.2cm]{img/1977/anhui_science/q3_fig1.png}
\end{problem}

\begin{solution}
\begin{enumerate}
\item 令二项式的第一项为 \(u=\sqrt[3]{x}=x^{\frac{1}{3}}\)，第二项为 \(v=-\frac{2}{3\sqrt[5]{x}}=-\frac{2}{3}x^{-\frac{1}{5}}\). 根据二项式定理，第 \(\left(r+1\right)\) 项为
\[
\begin{gathered}
T_{r+1}=\mathrm C_n^r u^{n-r}v^r
=\left(-1\right)^r\mathrm C_n^r\left(\frac{2}{3}\right)^r x^{\frac{n-r}{3}-\frac{r}{5}} \\
=\left(-1\right)^r\mathrm C_n^r\left(\frac{2}{3}\right)^r x^{\frac{5n-8r}{15}}
\end{gathered}
\]
其中 \(n\) 为自然数，\(0\leqslant r\leqslant n\)，且原式有意义时 \(x\neq0\).

\item 因为三角形 \(ABC\) 在 \(A\) 处为直角，且 \(M\) 是斜边 \(BC\) 的中点，所以直角三角形斜边中线定理给出 \(AM=BM=CM=\frac{1}{2}BC\). 从而
\[
8AM^2=8\left(\frac{BC}{2}\right)^2=2BC^2
\]
另一方面，由勾股定理，\(AB^2+AC^2=BC^2\)，于是
\[
AB^2+AC^2+BC^2=BC^2+BC^2=2BC^2=8AM^2
\]
这就证明了所给等式.
\end{enumerate}
\end{solution}

\begin{problem}
\begin{enumerate}
    \item 化简：\(\sqrt{5^{2\log_5\left( \lg x \right)}-2\lg x+1}\).
    \item 求下面数列的前 \(n\) 项和：\(\lg\frac{1}{3}+\lg\frac{10}{3^2}+\lg\frac{100}{3^4}+\lg\frac{1000}{3^8}+\dots\)
\end{enumerate}
\end{problem}

\begin{solution}
\begin{enumerate}
\item 因为 \(\lg x\) 是对数 \(\log_5\left(\lg x\right)\) 的真数，首先必须有 \(\lg x>0\)，即 \(x>1\). 在此条件下，利用 \(5^{\log_5 t}=t\)，有
\[
\sqrt{5^{2\log_5\left(\lg x\right)}-2\lg x+1}
=\sqrt{\left(\lg x\right)^2-2\lg x+1}
=\sqrt{\left(\lg x-1\right)^2}=\left|\lg x-1\right|
\]
当 \(x\geqslant10\) 时，\(\lg x-1\geqslant0\)，原式等于 \(\lg x-1=\lg\frac{x}{10}\)；当 \(1<x<10\) 时，\(\lg x-1<0\)，原式等于 \(1-\lg x=\lg\frac{10}{x}\). 因此
\[
\sqrt{5^{2\log_5\left(\lg x\right)}-2\lg x+1}
=\begin{cases}
\lg\dfrac{10}{x}, & 1<x<10,\\
\lg\dfrac{x}{10}, & x\geqslant10.
\end{cases}
\]

\item 第 \(k\) 项可写为 \(\lg\dfrac{10^{k-1}}{3^{2^{k-1}}}\)，所以前 \(n\) 项和为
\[
S_n=\sum_{k=1}^{n}\lg\frac{10^{k-1}}{3^{2^{k-1}}}
\]
由对数的商、幂运算公式，
\[
\begin{gathered}
S_n=\sum_{k=1}^{n}\left((k-1)-2^{k-1}\lg3\right)
=\frac{n\left(n-1\right)}{2}-\left(1+2+\cdots+2^{n-1}\right)\lg3 \\
=\frac{n\left(n-1\right)}{2}-\left(2^n-1\right)\lg3
=\frac{n\left(n-1\right)}{2}+\left(1-2^n\right)\lg3
\end{gathered}
\]
\end{enumerate}
\end{solution}

\begin{problem}
已知 \(\triangle ABC\) 的三个内角 \(A\)，\(B\)，\(C\) 成等差数列，且最大角与最小角的对边之比是 \(\left( \sqrt{3}+1 \right):2\)，试求此三角形三个内角的度数.
\end{problem}

\begin{solution}
设最大角、次大角、最小角分别为 \(U\)，\(V\)，\(W\)，并令其对边分别为 \(u\)，\(v\)，\(w\). 三角形内角成等差数列，所以 \(U+W=2V\). 又 \(U+V+W=180^\circ\)，从而 \(3V=180^\circ\)，即 \(V=60^\circ\)，且 \(U+W=120^\circ\). 题设的对边之比大于 \(1\)，故 \(U>W\)；结合角的大小次序可知 \(0^\circ<W<60^\circ<U<120^\circ\).

由正弦定理和题设边长比，
\[
\frac{\sin U}{\sin W}=\frac{u}{w}=\frac{\sqrt{3}+1}{2}
\]
因为 \(U=120^\circ-W\)，所以
\[
\frac{\sin U}{\sin W}
=\frac{\sin\left(120^\circ-W\right)}{\sin W}
=\sin120^\circ\cot W-\cos120^\circ
=\frac{\sqrt{3}}{2}\cot W+\frac{1}{2}
\]
于是
\[
\frac{\sqrt{3}}{2}\cot W+\frac{1}{2}=\frac{\sqrt{3}+1}{2}
\]
解得 \(\cot W=1\). 由于 \(0^\circ<W<60^\circ\)，故 \(W=45^\circ\)，进而 \(U=120^\circ-45^\circ=75^\circ\). 因此三角形的三个内角为 \(75^\circ\)，\(60^\circ\)，\(45^\circ\).
\end{solution}

\begin{problem}
设 \(P\left( 8,2 \right)\) 为平面直角坐标系上一点，\(y^2=x\) 是曲线，\(y-x-1=0\) 是一直线：

\begin{enumerate}
    \item 求过 \(P\) 点与直线 \(y-x-1=0\) 平行的直线方程.
    \item 求所求的直线与已知曲线交点的坐标.
    \item 作出它们的图象.
\end{enumerate}
\end{problem}

\begin{solution}
\begin{enumerate}
\item 已知直线可写为 \(y=x+1\)，斜率为 \(1\). 所求平行线过点 \(P\left(8,2\right)\)，所以由点斜式
\[
y-2=1\left(x-8\right)
\]
整理得所求直线方程为 \(y-x+6=0\).

\item 交点 \((x,y)\) 同时满足直线方程和曲线方程，即
\[
\begin{cases}
y-x+6=0,\\
y^2=x.
\end{cases}
\]
由第一式得 \(y=x-6\)，代入第二式，得
\[
(x-6)^2=x\quad\Longleftrightarrow\quad x^2-13x+36=0
\quad\Longleftrightarrow\quad (x-4)(x-9)=0
\]
所以 \(x=4\) 或 \(x=9\). 分别代入 \(y=x-6\)，得到 \((x,y)=\left(4,-2\right)\) 或 \(\left(9,3\right)\). 代回原方程可知两点均满足方程组，故交点为 \(\left(4,-2\right)\) 和 \(\left(9,3\right)\).

\item 曲线 \(y^2=x\) 是顶点为原点、开口向右的抛物线，所求直线为斜率为 \(1\) 且过 \(P\left(8,2\right)\) 的直线；它与抛物线交于 \(\left(4,-2\right)\) 和 \(\left(9,3\right)\). 图象如下：

\solutionfigure{\bitmapfigure[width=8.0cm]{img/1977/anhui_science/q6_solution_fig1.png}}
\end{enumerate}
\end{solution}

\begin{problem}
判断下列方程的轨迹，并画图象：\(25x^2-14xy+25y^2-288=0\).
\end{problem}

\begin{solution}
作坐标旋转变换，令坐标系逆时针旋转 \(\varphi\)，并记新坐标为 \((x',y')\). 坐标变换为
\(x=x'\cos\varphi-y'\sin\varphi\)，\(y=x'\sin\varphi+y'\cos\varphi\).
令 \(c=\cos\varphi\)，\(s=\sin\varphi\). 代入原方程并展开，得到
\[
\begin{gathered}
25(x'c-y's)^2-14(x'c-y's)(x's+y'c)+25(x's+y'c)^2-288=0 \\
\Longleftrightarrow\ (25-14sc)x'^2+(25+14sc)y'^2-14(c^2-s^2)x'y'-288=0
\end{gathered}
\]
为消去交叉项，取 \(c^2-s^2=0\)，可取 \(\varphi=\frac{\pi}{4}\). 此时 \(sc=\frac{1}{2}\)，方程化为
\[
18x'^2+32y'^2-288=0
\quad\Longleftrightarrow\quad
\frac{x'^2}{16}+\frac{y'^2}{9}=1
\]
这是以原点为中心、在新坐标轴方向上的椭圆，半长轴为 \(4\)，半短轴为 \(3\). 由于 \(x'\) 轴与原坐标系的直线 \(y=x\) 重合，\(y'\) 轴与直线 \(y=-x\) 重合，所以原坐标系中的图形是沿两条互相垂直的对角线方向伸长的椭圆，其图象如下：

\solutionfigure{\bitmapfigure[width=7.2cm]{img/1977/anhui_science/q7_solution_fig1.png}}
\end{solution}

\begin{problem}

\begin{enumerate}
    \item 用 \(0\) 到 \(9\) 这 \(10\) 个数字，可以排列成多少个含有两个重复数字的三位数？
    \item 用数学归纳法证明 \(x^n-y^n\) 能够被 \(x-y\) 整除（\(n\) 为自然数）.
\end{enumerate}
\end{problem}

\begin{solution}
\begin{enumerate}
\item 三位数的百位不能为 \(0\)，所以允许重复数字时，百位有 \(9\) 种选择，十位和个位各有 \(10\) 种选择，共有 \(9\times10\times10=900\) 个. 三个数字互不相同时，百位有 \(9\) 种选择，十位有除百位外的 \(9\) 种选择，个位有 \(8\) 种选择，共有 \(9\times9\times8=648\) 个. 三个数字完全相同时只能是 \(111\)，\(222\)，\(\dots\)，\(999\)，共有 \(9\) 个.

含有两个重复数字是指恰有一个数字出现两次，因此所求个数为
\[
900-648-9=243
\]

\item 先验证初始情形：当 \(n=1\) 时，\(x^1-y^1=x-y\)，显然能被 \(x-y\) 整除.

设当 \(n=k\) 时命题成立，即存在整式 \(q\)，使得 \(x^k-y^k=(x-y)q\). 则
\[
x^{k+1}-y^{k+1}
=x\left(x^k-y^k\right)+y^k\left(x-y\right)
=(x-y)\left(xq+y^k\right)
\]
因此 \(x^{k+1}-y^{k+1}\) 也能被 \(x-y\) 整除，命题从 \(k\) 推到了 \(k+1\). 根据数学归纳法，\(x^n-y^n\) 对所有自然数 \(n\geqslant1\) 都能被 \(x-y\) 整除；若将 \(0\) 也计入自然数，\(n=0\) 时差为 \(0\)，结论同样成立.
\end{enumerate}
\end{solution}


\section{附加题}

\begin{problem}
如图，在矩形 \(ABCD\) 中，\(A_1E\parallel AB\)，\(AA_1=A_1A_2=A_2A_3=A_3A_4=A_4D=a\)，\(AB_1=B_1B_2=B_2B=\sqrt{3}a\)，试用三种方法证明：\(\angle B_1A_1E+\angle B_2A_1E+\angle BDC=\frac{\pi}{2}\).
（注：三种方法是指每种方法各以代数，三角，几何知识为主来证明.）

\bitmapfigure[width=4.8cm]{img/1977/anhui_science/q9_fig1.png}
\end{problem}

\begin{solution}
法一：因为 \(ABCD\) 是矩形，\(A_1E\parallel AB\)，所以 \(\angle B_1A_1E\) 与直角三角形 \(AA_1B_1\) 在 \(B_1\) 处的锐角相等，故
\[
\tan\angle B_1A_1E=\tan\angle A_1B_1A=\frac{AA_1}{AB_1}=\frac{a}{\sqrt{3}a}=\frac{1}{\sqrt{3}}
\]
从而 \(\angle B_1A_1E=\frac{\pi}{6}\). 同理，
\[
\tan\angle B_2A_1E=\tan\angle A_1B_2A=\frac{AA_1}{AB_2}=\frac{a}{2\sqrt{3}a}=\frac{\sqrt{3}}{6}
\]
\[
\tan\angle BDC=\tan\angle DBA=\frac{AD}{AB}=\frac{5a}{3\sqrt{3}a}=\frac{5\sqrt{3}}{9}
\]
三个角均为锐角，且
\[
\tan\left( \angle B_2A_1E+\angle BDC \right)
=\frac{\frac{\sqrt{3}}{6}+\frac{5\sqrt{3}}{9}}{1-\frac{\sqrt{3}}{6}\cdot\frac{5\sqrt{3}}{9}}
=\sqrt{3}
\]
又 \(0<\angle B_2A_1E+\angle BDC<\pi\)，且上式的值为正，所以该和角为锐角；因此 \(\angle B_2A_1E+\angle BDC=\frac{\pi}{3}\). 因而
\[
\angle B_1A_1E+\angle B_2A_1E+\angle BDC
=\frac{\pi}{6}+\frac{\pi}{3}=\frac{\pi}{2}
\]

法二：延长 \(DA\) 到 \(A'\)，使 \(AA'=a\)，过 \(A'\) 作 \(A'B'\parallel AB\)，并令它与 \(BC\) 的延长线交于 \(B'\). 在 \(A'B'\) 上取点 \(B_2'\)，使 \(A'B_2'=2\sqrt{3}a\)，连接 \(DB_2'\) 和 \(B_2'B\)，如图所示：

\solutionfigure{\bitmapfigure[width=4.8cm]{img/1977/anhui_science/q9_solution_fig1.png}}

在直角三角形 \(B_1AA_1\) 和 \(DA'B_2'\) 中，\(DA'=6a\)，且
\(\frac{B_1A}{DA'}=\frac{\sqrt{3}a}{6a}=\frac{\sqrt{3}}{6}\)，\(\frac{AA_1}{A'B_2'}=\frac{a}{2\sqrt{3}a}=\frac{\sqrt{3}}{6}\).
两直角三角形的两组直角边对应成比例，且夹角均为直角，故由两边及其夹角判定 \(\triangle B_1AA_1\sim\triangle DA'B_2'\). 因此
\[
\angle B_1A_1E=\angle A_1B_1A=\angle A'DB_2'
\]

又在 \(\triangle DB_2'B\) 中，\(BD=\sqrt{\left(5a\right)^2+\left(3\sqrt{3}a\right)^2}\)，\(DB_2'=\sqrt{\left(6a\right)^2+\left(2\sqrt{3}a\right)^2}\)，\(B_2'B=\sqrt{\left(\sqrt{3}a\right)^2+a^2}\)，所以
\(BD^2=52a^2\)，\(DB_2'^2=48a^2\)，\(B_2'B^2=4a^2\).
于是 \(DB_2'^2+B_2'B^2=BD^2\)，由勾股定理的逆定理，\(\angle DB_2'B=\frac{\pi}{2}\).

在直角三角形 \(DB_2'B\) 和 \(B_2AA_1\) 中，
\(\frac{DB_2'}{B_2A}=\frac{\sqrt{48a^2}}{2\sqrt{3}a}=2\)，\(\frac{B_2'B}{AA_1}=\frac{2a}{a}=2\).
故这两个直角三角形相似，且顶点对应为 \(D\leftrightarrow B_2\)，\(B\leftrightarrow A_1\)，\(B_2'\leftrightarrow A\). 因而
\[
\angle B_2A_1E=\angle AB_2A_1=\angle B_2'DB
\]

射线 \(DB_2'\)，\(DB\) 位于射线 \(DA'\) 与 \(DC\) 之间，所以
\[
\angle B_1A_1E+\angle B_2A_1E+\angle BDC
=\angle A'DB_2'+\angle B_2'DB+\angle BDC
=\angle A'DC=\frac{\pi}{2}
\]

法三：在矩形 \(ABCD\) 所在平面内建立直角坐标系，以 \(A\) 为原点，\(AB\) 所在直线为 \(x\) 轴，\(AD\) 所在直线为 \(y\) 轴，并以 \(a\) 为长度单位. 由题设，相关点的坐标为 \(A_1\left(0,1\right)\)，\(B_1\left(\sqrt{3},0\right)\)，\(B_2\left(2\sqrt{3},0\right)\)，\(B\left(3\sqrt{3},0\right)\)，\(D\left(0,5\right)\). 于是直线 \(B_1A_1\)，\(B_2A_1\)，\(BD\) 的斜率分别为
\(k_1=\frac{1-0}{0-\sqrt{3}}=-\frac{1}{\sqrt{3}}=-\frac{\sqrt{3}}{3}\)，\(k_2=\frac{1-0}{0-2\sqrt{3}}=-\frac{1}{2\sqrt{3}}=-\frac{\sqrt{3}}{6}\)，\(k_3=\frac{5-0}{0-3\sqrt{3}}=-\frac{5}{3\sqrt{3}}=-\frac{5\sqrt{3}}{9}\).

因为 \(A_1E\) 与 \(x\) 轴平行，斜率为 \(0\). 两直线斜率为 \(m_1\)，\(m_2\) 时，夹角 \(\theta\) 满足 \(\tan\theta=\left|\frac{m_2-m_1}{1+m_1m_2}\right|\). 因此
\[
\tan\angle B_1A_1E
=\frac{0-\left( -\frac{\sqrt{3}}{3} \right)}{1+0\cdot\left( -\frac{\sqrt{3}}{3} \right)}
=\frac{\sqrt{3}}{3}
\]
故 \(\angle B_1A_1E=\frac{\pi}{6}\). 同理，

\[
\tan\angle B_2A_1E
=\frac{0-\left( -\frac{\sqrt{3}}{6} \right)}{1+0\cdot\left( -\frac{\sqrt{3}}{6} \right)}
=\frac{\sqrt{3}}{6}
\]

\[
\tan\angle BDC
=\frac{0-\left( -\frac{5\sqrt{3}}{9} \right)}{1+0\cdot\left( -\frac{5\sqrt{3}}{9} \right)}
=\frac{5\sqrt{3}}{9}
\]
又因 \(DC\) 与 \(x\) 轴平行，
\[
\tan\angle BDC=\frac{5\sqrt{3}}{9}
\]

代入正切的和角公式，
\[
\tan\left( \angle B_2A_1E+\angle BDC \right)
=\frac{\frac{\sqrt{3}}{6}+\frac{5\sqrt{3}}{9}}{1-\frac{\sqrt{3}}{6}\cdot\frac{5\sqrt{3}}{9}}
=\sqrt{3}
\]
两角之和属于 \((0,\pi)\)，而上式表明其正切为正，故两角之和为锐角；于是 \(\angle B_2A_1E+\angle BDC=\frac{\pi}{3}\). 从而
\[
\angle B_1A_1E+\angle B_2A_1E+\angle BDC
=\frac{\pi}{6}+\frac{\pi}{3}=\frac{\pi}{2}
\]
\end{solution}
